% !TeX program = lualatex
%
% Exemplo do pacote crpsp-revisao sobre a classe crpsp-timbrado: documento em
% revisão, com "USO INTERNO" no cabeçalho e a marca d'água "EM REVISÃO" atrás
% do texto. A arte (crpsp-timbrado-logo.pdf) não vem no repositório; ver o
% README. Sem a opção `marcadagua` só o cabeçalho muda: exemplo-revisao-simples.tex.
%
\DocumentMetadata{
  lang        = pt-BR,
  pdfversion  = 2.0,
  pdfstandard = ua-2,
  tagging     = on,
}
\documentclass{crpsp-timbrado}
\usepackage[marcadagua]{crpsp-revisao}

\hypersetup{
  pdftitle  = {Exemplo de documento em revisão},
  pdfauthor = {CRP SP},
}

\begin{document}

\noindent São Paulo, 19 de outubro de 2026.

\section{Objeto}

Este exemplo mostra o pacote \texttt{crpsp-revisao} com a marca d'água ligada.
O texto abaixo repete a estrutura do exemplo do papel timbrado, para conferir a
legibilidade sobre o fundo vermelho transparente.

\subsection{Cabeçalho}

À direita do logo aparece \textbf{USO INTERNO}, e a indicação de página fica
logo abaixo, alinhada à margem direita do bloco de texto.

\subsubsection{Marca d'água}

A marca é centrada no bloco de texto, e não na folha, girada em 45 graus e
composta atrás de todo o conteúdo.

\paragraph{Quarto nível}

Texto sob o título de nível 4, sobre a marca d'água.

\subparagraph{Quinto nível}

Texto sob o título de nível 5.

\section{Encaminhamentos}

São os seguintes os encaminhamentos:

\begin{enumerate}
  \item avaliar a minuta do documento;
  \item solicitar parecer da Gerência Jurídica;
  \item publicar a versão final em formato acessível.
\end{enumerate}

\subsection{Prazos}

\begin{tabular}{@{}ll@{}}
  \textbf{Etapa}  & \textbf{Prazo}    \\
  Avaliação       & 30 de outubro     \\
  Parecer         & 13 de novembro    \\
  Publicação      & 27 de novembro    \\
\end{tabular}

\section{Texto longo}

Parágrafo de enchimento para forçar uma segunda página e exercitar a indicação
de número de página empilhada sob o rótulo. A psicóloga inscrita no Conselho
Regional exerce função pública e está sujeita à fiscalização ética; o órgão
zela pela qualidade técnica, pela justiça e pela inclusão das pessoas.

\count255=0
\loop\ifnum\count255<12
  \advance\count255 by 1
  Três princípios sustentam a atuação: o compromisso com os direitos humanos, a
  defesa da ciência e a construção coletiva. Não há exercício profissional sem
  responsabilidade, nem avaliação séria sem informação clara.\par
\repeat

\end{document}
